% =============================================================================
% 九章算術·卷四至卷六 — 文言文↔白话文双语对照版
% Jiuzhang Suanshu (Nine Chapters) Volumes 4-6: Classical↔Modern Chinese
% =============================================================================
\documentclass[11pt,a4paper]{article}

% -------- Core Packages --------
\usepackage{fontspec}
\usepackage{xeCJK}
\usepackage{polyglossia}
\setmainlanguage{english}

% -------- CJK Font Configuration --------
\setCJKmainfont[AutoFakeBold=true,AutoFakeSlant=true]{Noto Sans CJK SC}
\setCJKsansfont[AutoFakeBold=true,AutoFakeSlant=true]{Noto Sans CJK SC}

% -------- Line breaking --------
\XeTeXlinebreaklocale "zh"
\XeTeXlinebreakskip = 0pt plus 1pt
\tolerance=1200
\emergencystretch=3em
\hyphenpenalty=5000
\exhyphenpenalty=5000

\usepackage{amsmath,amssymb,amsthm}
\usepackage{geometry}
\usepackage{graphicx}
\usepackage{xcolor}
\usepackage{titlesec}
\usepackage{fancyhdr}
\usepackage{enumitem}
\usepackage{array,booktabs,longtable,multirow}
\usepackage{hyperref}
\usepackage{tocloft}
\usepackage{indentfirst}
\usepackage{ragged2e}

% -------- Parallel Columns --------
\usepackage{paracol}
\colseprulecolor{rulecolor}
\setlength{\columnseprule}{0.4pt}

% -------- Page Geometry --------
\geometry{
  paperwidth=210mm,paperheight=297mm,
  left=18mm,right=18mm,top=25mm,bottom=25mm,
  headheight=14pt,footskip=30pt
}

% -------- Font Configuration --------
\setmainfont{Liberation Serif}
\setsansfont{Liberation Sans}
\setmonofont[Scale=MatchLowercase]{Liberation Mono}

% -------- Colors --------
\definecolor{titlecolor}{RGB}{45,55,72}
\definecolor{sectioncolor}{RGB}{55,65,81}
\definecolor{volcolor}{RGB}{120,53,15}
\definecolor{commentarycolor}{RGB}{80,60,40}
\definecolor{rulecolor}{RGB}{180,160,140}
\definecolor{problemcolor}{RGB}{25,55,95}
\definecolor{answercolor}{RGB}{35,85,55}
\definecolor{methodcolor}{RGB}{85,35,35}

% -------- Section Formatting --------
\titleformat{\section}
  {\normalfont\Large\bfseries\color{titlecolor}}
  {\thesection}{1em}{}
\titleformat{\subsection}
  {\normalfont\large\bfseries\color{sectioncolor}}
  {\thesubsection}{1em}{}
\titleformat{\subsubsection}
  {\normalfont\normalsize\bfseries\color{sectioncolor}}
  {\thesubsubsection}{1em}{}

% -------- Header/Footer --------
\pagestyle{fancy}
\fancyhf{}
\fancyhead[L]{\small\CJKfamily{rm}九章算術·卷四至卷六\ —\ 文言文白话文对照}
\fancyhead[R]{\small\thepage}
\renewcommand{\headrulewidth}{0.4pt}

% -------- Custom Environments --------
\usepackage{tcolorbox}
\tcbuselibrary{skins,breakable}

% Problem environment (bilingual wrapper)
\newtcolorbox{problembox}[1][]{
  colback=blue!3!white,colframe=problemcolor!60!black,
  fonttitle=\bfseries,title={#1},
  breakable,enhanced,boxrule=0.5pt,arc=2pt,
  before upper={\RaggedRight}
}

% Answer environment
\newtcolorbox{answerbox}[1][]{
  colback=green!3!white,colframe=answercolor!60!black,
  fonttitle=\bfseries,title={#1},
  breakable,enhanced,boxrule=0.5pt,arc=2pt,
  before upper={\RaggedRight}
}

% Method environment
\newtcolorbox{methodbox}[1][]{
  colback=red!3!white,colframe=methodcolor!60!black,
  fonttitle=\bfseries,title={#1},
  breakable,enhanced,boxrule=0.5pt,arc=2pt,
  before upper={\RaggedRight}
}

% Commentary environment
\newtcolorbox{commentarybox}[1][]{
  colback=yellow!4!white,colframe=commentarycolor!60!black,
  fonttitle=\bfseries\itshape,title={#1},
  breakable,enhanced,boxrule=0.4pt,arc=2pt,
  before upper={\RaggedRight}
}

% Volume header
\newtcolorbox{volumeheader}[2][]{
  colback=volcolor!8!white,colframe=volcolor!70!black,
  fonttitle=\bfseries\Large,title={#2},#1,
  breakable,enhanced,boxrule=1pt,arc=3pt
}

% -------- Custom Commands --------
\newcommand{\longline}{\noindent\rule{\textwidth}{0.4pt}}
\newcommand{\shortline}{\noindent\rule{0.5\textwidth}{0.4pt}\par}

% Bilingual section headers
\newcommand{\bilingualsection}[2]{%
  \section{#1\hfill\normalfont\large\color{sectioncolor}#2}%
}
\newcommand{\bilingualsubsection}[2]{%
  \subsection{#1\hfill\normalfont\normalsize\color{sectioncolor}#2}%
}

% Hyperref
\hypersetup{
  colorlinks=true,linkcolor=sectioncolor,citecolor=sectioncolor,
  urlcolor=sectioncolor,pdfencoding=auto
}

\setcounter{tocdepth}{2}
\setcounter{secnumdepth}{2}

% =============================================================================
\newcommand{\vs}{\vspace{0.5em}}
\begin{document}

% ------------------------------------------------------------------------------
% TITLE PAGE
% ------------------------------------------------------------------------------
\begin{titlepage}
\begin{center}
\vspace*{2cm}

{\Huge\bfseries\color{titlecolor} 九章算術}\\[0.8em]
{\LARGE\color{sectioncolor} 卷四至卷六}\\[1.5em]
{\Large\color{volcolor} 文言文 \raisebox{0.1ex}{$\leftrightarrow$} 白话文对照版}\\[2.5em]

\shortline

\vspace{1.5cm}

{\large\bfseries 第四卷：少廣}\\[0.3em]
{\normalsize 开方求积、球体积}\\[1em]
{\large\bfseries 第五卷：商功}\\[0.3em]
{\normalsize 城墙、堤坝、沟渠体积}\\[1em]
{\large\bfseries 第六卷：均输}\\[0.3em]
{\normalsize 加权分配、行程问题}\\[2em]

\shortline

\vspace{1.5cm}

{\normalsize （晋）刘徽 注 \quad | \quad （唐）李淳风 注释}\\[2em]

{\small 据《钦定四库全书》子部·天文算法类}\\[0.5em]
{\small 原文（左栏）：文言文 \quad 译文（右栏）：现代汉语白话文}

\vfill

{\small 排版：XeLaTeX + Noto Sans CJK SC}

\end{center}
\end{titlepage}

% ------------------------------------------------------------------------------
% TABLE OF CONTENTS
% ------------------------------------------------------------------------------
\tableofcontents
\newpage

% ------------------------------------------------------------------------------
% INTRODUCTION
% ------------------------------------------------------------------------------
\bilingualsection{简介}{Introduction}

\begin{paracol}{2}

\textbf{原文：}\\
《九章算术》是中国古代最重要的数学经典，成书于汉代（公元前202年—公元220年）。全书分为九章，每章针对一类实际数学问题。

刘徽（约225—295年）于263年所注，是中国古代数学最伟大的成就之一。他不仅解释了算法，更以"出入相补"之法给出严格证明，并以内接正3072边形推算圆周率近似值。

李淳风（602—670年）奉唐诏主持注释，进一步阐明与校正原文。

\switchcolumn

\textbf{白话文：}\\
《九章算术》是中国古代最重要的数学经典，编纂于汉代（公元前202年—公元220年）。全书共分九章（卷），每一章都针对一类实际数学问题。

刘徽（约225—295年）在263年完成的注释，是数学史上最杰出的著作之一。刘徽不仅解释了各种算法，更运用"以盈补虚"（出入相补）的方法给出了严格的几何证明，并以内接正3072边形计算出圆周率的精确近似值。

李淳风（602—670年）受唐朝皇帝诏令，主持了对《九章算术》的进一步注释和校正工作。

\end{paracol}

\vspace{1cm}

\begin{center}
\begin{longtable}{|c|l|p{4.5cm}|}
\hline
\textbf{卷} & \textbf{篇名} & \textbf{内容} \\
\hline
\endfirsthead
\hline
\textbf{卷} & \textbf{篇名} & \textbf{内容} \\
\hline
\endhead
四 & 少广 & 开方（平方根、立方根）、已知面积/体积求边长、球体积 \\
五 & 商功 & 土方工程体积（城墙、堤坝、沟渠、粮仓等） \\
六 & 均输 & 公平赋税、加权分配、行程与工作效率问题 \\
\hline
\end{longtable}
\end{center}

\vspace{0.5cm}

\noindent
\textbf{体例说明：}每道题遵循古典结构：
\begin{itemize}[leftmargin=2em]
  \item \textbf{问}（\textcolor{problemcolor}{Problem}）：题目陈述——"今有……"
  \item \textbf{答}（\textcolor{answercolor}{Answer}）：答案——"答曰……"
  \item \textbf{术}（\textcolor{methodcolor}{Method}）：算法——"术曰……"
  \item \textbf{注}（\textcolor{commentarycolor}{Commentary}）：刘徽注释，解释推理原理
\end{itemize}

\newpage

% =============================================================================
\bilingualsection{卷四：少广}{Vol.\ 4: Shaoguang (Lesser Breadth)}
% =============================================================================

\begin{volumeheader}{\CJKfamily{rm}九章算術卷四 · 少广}
少广以御积幂方圆\\[0.3em]
\small 白话文：少广章用于处理面积、幂次、方圆等计算
\end{volumeheader}

\vspace{0.5cm}

\begin{paracol}{2}
少广章研究已知矩形面积和一个边长，求另一边长的问题。引申到开方术（平方根、立方根），以及著名的球体积计算问题。"少广"之意，即在已知大面积/体积的情况下，求其"少"（较小）的边长尺寸。

\switchcolumn

少广章研究的问题是：已知一个矩形的面积和其中一条边的长度，如何求出另一条边的长度。这一方法进一步延伸到开平方、开立方等开方运算，以及著名的球体积计算难题。"少广"这个标题的含义是：在已知较大面积或体积的情况下，反推出其中较小的那条边长。

\end{paracol}

\vspace{0.5cm}

% ------------------------------------------------------------------------------
\bilingualsubsection{少广术：开方法}{Opening Method: The Shaoguang Algorithm}
% ------------------------------------------------------------------------------

\begin{paracol}{2}

\begin{methodbox}[\CJKfamily{rm}少广术 · 原文]
置全步及分母子，以最下分母遍乘诸分子及全步。各以其母除其子，置之於左。命通分，内子。又以分母遍乘诸分子，及已通者，皆通而同之，并之为法。置所求步数，以全步积分乘之，为实。实如法而一，得从步。
\end{methodbox}

\switchcolumn

\begin{methodbox}[白话文：少广求边长算法]
列出整数步长和各个分数的分子、分母。用最大的分母去遍乘所有的分子和整数步长。然后各自用分母去除分子，将结果放在左边。进行通分运算，将分子纳入。再用分母遍乘各分子以及已经通分的数，使所有分数都通分为同分母，相加作为除数（法）。将所要求的步长数，用整数步长的总积分去乘，作为被除数（实）。被除数除以除数，便得到所求的长度步数。
\end{methodbox}

\end{paracol}

\vspace{0.3cm}

\begin{paracol}{2}

\begin{commentarybox}[\CJKfamily{rm}刘徽注 · 原文]
此以田广为法，以亩积步为实。术曰置全步及分母子者，通分而齐其子也。齐其子，则母同其母。以母乘全步者，通其分也。以母遍乘诸分子，各以其母除其子，置之於左。命通分内子者，通而同之也。
\end{commentarybox}

\switchcolumn

\begin{commentarybox}[白话文：刘徽注释]
这是用田地的宽度作为除数（法），用一亩的步数面积作为被除数（实）。"术曰：置全步及分母子"的意思是要通分，使各分数的分子统一。统一了分子，就要让各分数的分母也相同。用分母去乘整数步长，是为了给整数部分也通分。用分母遍乘各分子，再各自用分母除分子，放在左边。"命通分内子"的意思就是将各分数通分为同分母后再进行运算。
\end{commentarybox}

\end{paracol}

\vspace{0.5cm}
\longline
\vspace{0.3cm}

% ------------------------------------------------------------------------------
\bilingualsubsection{题一：已知面积带分数求边长}{Problem 1: Finding the Side from Area with Fractional Parts}
% ------------------------------------------------------------------------------

\begin{paracol}{2}

\begin{problembox}[\CJKfamily{rm}少广题一 · 原文]
\textbf{问：}今有田，广一步半。求田一亩，问从几何？
\end{problembox}

\switchcolumn

\begin{problembox}[白话文：少广第一题]
\textbf{题目：}现有一块田地，宽度为一步半（即1.5步）。如果要使这块田的恰好为一亩，问它的长度应该是多少？
\end{problembox}

\end{paracol}

\vspace{0.2cm}

\begin{paracol}{2}

\begin{answerbox}[\CJKfamily{rm}答 · 原文]
答曰：一百六十步。
\end{answerbox}

\switchcolumn

\begin{answerbox}[白话文：答案]
答案是：160步。
\end{answerbox}

\end{paracol}

\vspace{0.2cm}

\begin{paracol}{2}

\begin{methodbox}[\CJKfamily{rm}术 · 原文]
术曰：下有半，是二分之一。以一为二，半为一，并之得三，为法。置田二百四十步，亦以一为二乘之，为实。实如法得从步。
\end{methodbox}

\switchcolumn

\begin{methodbox}[白话文：解题方法]
算法说：下面有"半"，就是二分之一。将一化为二，半化为一，相加得三，作为除数（法）。将一亩对应的240步，也将一化为二来乘，作为被除数（实）。被除数除以除数，就得到所求的长度步数。
\end{methodbox}

\end{paracol}

\vspace{0.2cm}

\begin{paracol}{2}

\begin{commentarybox}[\CJKfamily{rm}刘徽注 · 原文]
此以广为法，以亩积步为实。术曰下有半是二分之一者，广一步半，其半即二分之一也。以一为二者，通分也。半为一者，二分之一即一也。并之得三为法者，分母二通全步一为二，分子一，共三也。
\end{commentarybox}

\switchcolumn

\begin{commentarybox}[白话文：刘徽注释]
这里用田地的宽度作为除数（法），用一亩的积步数作为被除数（实）。"术曰：下有半是二分之一"的意思是：宽度为一步半，其中"半"就是二分之一。"以一为二"是通分（化为以2为分母）。"半为一"是说二分之一在这种通分后就是一。"并之得三为法"的意思是：分母2将整数1通分为2，加上分子1，一共是3。
\end{commentarybox}

\end{paracol}

\vspace{0.5cm}
\longline
\vspace{0.3cm}

% ------------------------------------------------------------------------------
\bilingualsubsection{开方术：平方根求法}{Square Root Extraction}
% ------------------------------------------------------------------------------

\begin{paracol}{2}

\begin{problembox}[\CJKfamily{rm}开方 · 原文]
\textbf{问：}今有积五万五千二百二十五步。问为方几何？
\end{problembox}

\switchcolumn

\begin{problembox}[白话文：开平方问题]
\textbf{题目：}现有一个面积为五万五千二百二十五平方步。问这个正方形的边长是多少？
\end{problembox}

\end{paracol}

\vspace{0.2cm}

\begin{paracol}{2}

\begin{answerbox}[\CJKfamily{rm}答 · 原文]
答曰：二百三十五步。
\end{answerbox}

\switchcolumn

\begin{answerbox}[白话文：答案]
答案是：235步。
\end{answerbox}

\end{paracol}

\vspace{0.2cm}

\begin{paracol}{2}

\begin{methodbox}[\CJKfamily{rm}开方术 · 原文]
术曰：置积为实。借一算，步之，超二等。议所得，以一乘所借一算为法，而以除。除已，倍法为定法。其复除，折法而下。复置借算，步之如初。以复议一乘之，所得副，以加定法，以除。以所得副从定法。复除，折下如前。
\end{methodbox}

\switchcolumn

\begin{methodbox}[白话文：开平方算法]
算法说：将被开方数（面积）作为被除数（实）。借一个算筹用于定位，每步移动时跳过两位（即按百位递进）。商得第一位数后，用它乘借算作为除数（法），进行除法运算。除完后，将除数加倍作为固定的除数（定法）。继续除下一位时，将定数折下移。重新放置借算，像最初那样定位。再用商得的数乘它，所得结果作为副数，加到定法上，然后除。将副数并入定法。继续这样运算，依次折下求每一位。
\end{methodbox}

\end{paracol}

\vspace{0.2cm}

\begin{paracol}{2}

\begin{commentarybox}[\CJKfamily{rm}刘徽注 · 原文]
开方术者，求方幂之一面也。凡幂，方圆之本也。方幂者，方之积也。开方者，求方之面数。言百之面十，万之面百。术曰置积为实者，积即方幂也。借一算步之者，假以定位也。超二等者，方十自乘，幂一百；方百自乘，幂一万。故超两位也。
\end{commentarybox}

\switchcolumn

\begin{commentarybox}[白话文：刘徽论开方术]
开方术的目的，是求出正方形面积的一条边长。凡是幂（面积），都是方圆计算的基础。方幂就是正方形的面积。开方就是求正方形的边长。例如100的平方根是10，10000的平方根是100。"术曰：置积为实"中，被开方数就是方幂（面积）。"借一算"是用一个算筹来定位。"超二等"是因为10自乘得100，100自乘得10000，所以每进一位要跳两位。
\end{commentarybox}

\end{paracol}

\vspace{0.5cm}
\longline
\vspace{0.3cm}

% ------------------------------------------------------------------------------
\bilingualsubsection{立圆：球体积问题}{The Sphere Volume Problem}
% ------------------------------------------------------------------------------

这是《九章算术》中最著名的问题之一，刘徽对此进行了深入分析，为后世祖冲之、祖暅父子的工作奠定了基础。

\begin{paracol}{2}

\begin{problembox}[\CJKfamily{rm}立圆 · 原文]
\textbf{问：}今有立圆，径四尺。问为积几何？
\end{problembox}

\switchcolumn

\begin{problembox}[白话文：球体积问题]
\textbf{题目：}现有一个实心球体，直径为4尺。问它的体积是多少？
\end{problembox}

\end{paracol}

\vspace{0.2cm}

\begin{paracol}{2}

\begin{answerbox}[\CJKfamily{rm}答 · 原文]
答曰：九尺五寸八分寸之一。
\end{answerbox}

\switchcolumn

\begin{answerbox}[白话文：答案]
答案是：9尺5寸又1/8寸。（这是用古法公式 $V = \frac{9}{16}d^3$ 计算的结果。）
\end{answerbox}

\end{paracol}

\vspace{0.2cm}

\begin{paracol}{2}

\begin{methodbox}[\CJKfamily{rm}术 · 原文]
术曰：圆径自乘，九之，十六而一。开立圆术：置圆径四尺，自乘，得一十六尺。九之，得一百四十四尺。十六而一，得九尺。余十六分寸之八，约之为八分寸之一。并之，为积。
\end{methodbox}

\switchcolumn

\begin{methodbox}[白话文：球体积算法]
算法说：将球直径自乘，再乘以9，除以16。开立圆术的具体步骤是：球直径为4尺，自乘得16尺。乘以9得144尺。除以16得9尺。余下十六分之八尺，约分为八分之一尺。合并起来，就是球的体积。
\end{methodbox}

\end{paracol}

\vspace{0.2cm}

\begin{paracol}{2}

\begin{commentarybox}[\CJKfamily{rm}刘徽注 · 原文]
按此术，圆径自乘者，先得圆幂。九之十六而一者，以圆幂为方幂，九之十六而一，即圆积也。然此术以圆径为方径，误也。何以验之？取立方棋八枚，皆令立方一寸，积之为立方二寸。规之为圆，径二寸。又横规之，径亦二寸。然后并立圆棋二枚，规之为圆，径二寸。又横规之，径亦二寸。是为圆积也。

徽术曰：立方径自乘，又以圆周率三乘之，十二而一，得圆积。此圆幂之率也。又以此圆积乘高，得圆柱积。立圆者，圆柱之半也。故又以圆周率三乘之，十二而一，得立圆积。此术未得其理。

牟合方盖者，方率也。丸居牟合方盖，则圆率也。按合盖者，方率也。其中则圆周率也。方圆相缠，周径相直。若设方率为四，圆周率为三。以圆周率乘方幂，开方除之，即径。以方率乘圆幂，开方除之，即周。此圆方之率也。
\end{commentarybox}

\switchcolumn

\begin{commentarybox}[白话文：刘徽球体积注释]
按照这个算法，先将圆直径自乘得到圆面积（圆幂）。再九之十六而一，是把圆面积当作方面积来用九比十六的比率换算，这就是球体积。但这个算法把圆的直径当成了正方形的边长，这是错误的。如何验证呢？取八枚立方棋，每枚边长一寸，合在一起组成边长二寸的立方体。用圆规画圆，直径二寸。再横向画圆，直径也是二寸。然后将两枚立圆棋合在一起画圆，直径二寸。再横向画圆，直径也是二寸。这才是圆积。

刘徽自己的算法说：将立方体边长自乘，再用圆周率3去乘，除以12，得到圆面积。这是圆面积的比率。再用这圆面积乘高，得到圆柱体积。球体可以看作是圆柱的一半，所以再用圆周率3去乘，除以12，得到球体积。但这个算法也还没能得到正确的原理。

"牟合方盖"（两个垂直圆柱的交合体）是方率，球体包在牟合方盖里面，就是圆率。牟合方盖代表方率，里面内切的球体代表圆率。方圆相互缠绕，周长和直径相互对应。如果设方率为4，圆周率为3。用圆周率乘方面积再开方就是直径。用方率乘圆面积再开方就是周长。这就是圆和方的比率关系。
\end{commentarybox}

\end{paracol}

\textit{注释说明：}刘徽在此处指出了古法 $V = \frac{9}{16}d^3$（暗含 $\pi = 3$ 的特定配置）的错误。他引入了\textbf{牟合方盖}（mouhe fangai，两个垂直圆柱的交合体）这一关键概念，正确地指出球体内切于此立体之中，但未能完成最终计算。这一问题后来在5世纪由\textbf{祖暅}（祖冲之之子）解决，他证明了牟合方盖的体积是其外切立方体的 $\frac{2}{3}$，从而导出正确的球体积公式 $V = \frac{\pi}{6}d^3$。

\newpage

% =============================================================================
\bilingualsection{卷五：商功}{Vol.\ 5: Shanggong (Construction Consultations)}
% =============================================================================

\begin{volumeheader}{\CJKfamily{rm}九章算術卷五 · 商功}
商功以御功程积实\\[0.3em]
\small 白话文：商功章用于处理工程土方量与劳动力计算
\end{volumeheader}

\vspace{0.5cm}

\begin{paracol}{2}
商功章研究各种工程建筑和土方计算中的体积问题：城墙、堤坝、沟渠、运河、粮仓等，以及丰富的几何体体积。刘徽在此章的注释尤为详尽，运用分割法给出几何证明。

\switchcolumn

商功章研究各种工程建筑和土方计算中的体积问题，包括城墙、堤坝、沟渠、运河、粮仓等工程体的体积计算，以及各类几何体的体积公式。刘徽在这一章的注释特别详尽，运用几何分割和拼补的方法给出了各种体积公式的严格证明。

\end{paracol}

\vspace{0.5cm}

% ------------------------------------------------------------------------------
\bilingualsubsection{穿地率：土方分类}{Classification of Earthwork Solids}
% ------------------------------------------------------------------------------

\begin{paracol}{2}

\begin{methodbox}[\CJKfamily{rm}商功术 · 原文]
穿地四，为壤五，为坚三。墟谓穿坑，此皆其常率。
\end{methodbox}

\switchcolumn

\begin{methodbox}[白话文：土方换算率]
挖出的松散土方（穿地）4份，堆起来会变成5份（壤），夯实后变为3份（坚）。"墟"就是指挖出的坑——这些都是标准的换算率。
\end{methodbox}

\end{paracol}

\vspace{0.2cm}

\begin{paracol}{2}

\begin{commentarybox}[\CJKfamily{rm}刘徽注 · 原文]
以穿地求壤，五之；求坚，三之。皆四而一。以壤求穿，四之；求坚，三之。皆五而一。以坚求穿，四之；求壤，五之。皆三而一。此术皆其常率也。穿地四为壤五者，掘地四尺，出壤五尺，壤虚而坚实也。
\end{commentarybox}

\switchcolumn

\begin{commentarybox}[白话文：刘徽注释]
由挖出的土方（穿地）求松散土方（壤），乘以5；求夯实土方（坚），乘以3。都除以4。由松散土方求挖出土方，乘以4；求夯实土方，乘以3。都除以5。由夯实土方求挖出土方，乘以4；求松散土方，乘以5。都除以3。这些都是标准的换算率。穿地4份变为壤5份，是因为掘地4尺，挖出来的土堆起来会变成5尺，因为松散的土蓬松而夯实的土密实。
\end{commentarybox}

\end{paracol}

\vspace{0.5cm}
\longline
\vspace{0.3cm}

% ------------------------------------------------------------------------------
\bilingualsubsection{城垣问题：城墙体积}{Problem: Volume of a City Wall}
% ------------------------------------------------------------------------------

\begin{paracol}{2}

\begin{problembox}[\CJKfamily{rm}城垣 · 原文]
\textbf{问：}今有城，下广四丈，上广二丈，高五丈，袤一百二十六丈。问积几何？
\end{problembox}

\switchcolumn

\begin{problembox}[白话文：城墙体积问题]
\textbf{题目：}现有一座城墙，底部宽4丈，顶部宽2丈，高5丈，长126丈。问它的体积是多少？
\end{problembox}

\end{paracol}

\vspace{0.2cm}

\begin{paracol}{2}

\begin{answerbox}[\CJKfamily{rm}答 · 原文]
答曰：一百八十九万尺。
\end{answerbox}

\switchcolumn

\begin{answerbox}[白话文：答案]
答案是：1,890,000立方尺。
\end{answerbox}

\end{paracol}

\vspace{0.2cm}

\begin{paracol}{2}

\begin{methodbox}[\CJKfamily{rm}术 · 原文]
术曰：并上下广而半之，以高若深乘之，又以袤乘之，即积尺。
\end{methodbox}

\switchcolumn

\begin{methodbox}[白话文：城墙体积算法]
算法说：将上宽和下宽相加后取半，乘以高度（或深度），再乘以长度，就得到体积（立方尺）。
\end{methodbox}

\end{paracol}

\vspace{0.2cm}

\begin{paracol}{2}

\begin{commentarybox}[\CJKfamily{rm}刘徽注 · 原文]
按此术，并上下广而半之者，以盈补虚，得中平之广。以高乘之，得截面为直棱之积。又以袤乘之者，广袤相乘为一截之积，以高乘之为积。
\end{commentarybox}

\switchcolumn

\begin{commentarybox}[白话文：刘徽注释]
按照这个算法，将上宽和下宽相加取半，是通过"以盈补虚"（用多余的部分填补不足的部分）的原理，得到中间平均的宽度。用高度去乘，就得到横截面的面积（相当于一个直棱柱的截面积）。再用长度去乘，就是横截面面积乘以长度得到总体积。
\end{commentarybox}

\end{paracol}

\vspace{0.5cm}
\longline
\vspace{0.3cm}

% ------------------------------------------------------------------------------
\bilingualsubsection{堑堵：沟渠体积}{Problem: Volume of a Canal/Trench}
% ------------------------------------------------------------------------------

\begin{paracol}{2}

\begin{problembox}[\CJKfamily{rm}堑堵 · 原文]
\textbf{问：}今有堑，上广一丈六尺，下广一丈，深六尺，袤五丈七尺。问积几何？
\end{problembox}

\switchcolumn

\begin{problembox}[白话文：沟渠体积问题]
\textbf{题目：}现有一条沟渠，上口宽1丈6尺，下底宽1丈，深6尺，长5丈7尺。问它的体积是多少？
\end{problembox}

\end{paracol}

\vspace{0.2cm}

\begin{paracol}{2}

\begin{answerbox}[\CJKfamily{rm}答 · 原文]
答曰：积七千一百一十二尺。
\end{answerbox}

\switchcolumn

\begin{answerbox}[白话文：答案]
答案是：体积为7,112立方尺。
\end{answerbox}

\end{paracol}

\vspace{0.2cm}

\begin{paracol}{2}

\begin{methodbox}[\CJKfamily{rm}术 · 原文]
术曰：并上下广，半之，以高若深乘之，又以袤乘之，即积尺。
\end{methodbox}

\switchcolumn

\begin{methodbox}[白话文：沟渠体积算法]
算法说：将上宽和下宽相加取半，乘以高度或深度，再乘以长度，就得到体积（立方尺）。
\end{methodbox}

\end{paracol}

\vspace{0.5cm}
\longline
\vspace{0.3cm}

% ------------------------------------------------------------------------------
\bilingualsubsection{阳马与鳖臑：基本几何体}{The Yangma and Bienao}
% ------------------------------------------------------------------------------

阳马和鳖臑是中国古代数学中最重要的几何体之一，通过对长方体的分割得到。

\begin{paracol}{2}

\begin{methodbox}[\CJKfamily{rm}阳马术 · 原文]
术曰：广袤相乘，以高乘之，三而一。

\vspace{0.3em}
此求\textbf{阳马}（直角方锥）之积：$V = \frac{1}{3} \times$ 宽 $\times$ 长 $\times$ 高。
\end{methodbox}

\switchcolumn

\begin{methodbox}[白话文：阳马（直角方锥）体积算法]
算法说：将宽度与长度相乘，再乘以高度，除以3。

\vspace{0.3em}
这就是求\textbf{阳马}的体积公式：$V = \frac{1}{3} \times$ 宽 $\times$ 长 $\times$ 高。
\end{methodbox}

\end{paracol}

\vspace{0.2cm}

\begin{paracol}{2}

\begin{commentarybox}[\CJKfamily{rm}刘徽注 · 原文]
按此术，阳马者，隅角之椎也。解立方得两堑堵，解堑堵得一阳马、一鳖臑。阳马居二，鳖臑居一，不易之率也。故立方三而一，即阳马之积也。

合二堑堵成一阳马者，试令广袤各六尺，高六尺。中分其高，令广袤各三尺。其高六尺，中分为二，各三尺。上堑堵广袤各三尺，高亦三尺；下堑堵亦然。合二堑堵，广袤各六尺，高六尺，是为一阳马也。
\end{commentarybox}

\switchcolumn

\begin{commentarybox}[白话文：刘徽论阳马]
按照这个算法，阳马是位于角落的锥体。将一个长方体分解得到两个堑堵（楔形），再将每个堑堵分解得到一个阳马和一个鳖臑。阳马占2份，鳖臑占1份，这是不变的比率。所以长方体体积除以3就是阳马的体积。

两个堑堵合成一个阳马的证明：假设宽和长都是6尺，高也是6尺。将高平分，使宽和长都变为3尺。高6尺平分为二，各3尺。上面的堑堵宽、长各3尺，高也是3尺；下面的堑堵也一样。两个堑堵合在一起，宽、长各6尺，高6尺，就是一个阳马。
\end{commentarybox}

\end{paracol}

\vspace{0.3cm}

\begin{paracol}{2}

\begin{methodbox}[\CJKfamily{rm}鳖臑术 · 原文]
术曰：广袤相乘，以高乘之，六而一。

\vspace{0.3em}
此求\textbf{鳖臑}（四面体）之积：$V = \frac{1}{6} \times a \times b \times c$。
\end{methodbox}

\switchcolumn

\begin{methodbox}[白话文：鳖臑（四面体）体积算法]
算法说：将宽度与长度相乘，再乘以高度，除以6。

\vspace{0.3em}
这就是求\textbf{鳖臑}的体积公式：$V = \frac{1}{6} \times a \times b \times c$。
\end{methodbox}

\end{paracol}

\vspace{0.2cm}

\begin{paracol}{2}

\begin{commentarybox}[\CJKfamily{rm}刘徽注 · 原文]
鳖臑者，推之明理也。按阳马之积，三分立方之一。鳖臑者，又半阳马也。故六而一。凡立方，高一尺，广袤各一尺，积一尺。阳马广袤各一尺，高一尺，积三分之一尺。鳖臑三边各一尺，积六分之一尺。
\end{commentarybox}

\switchcolumn

\begin{commentarybox}[白话文：刘徽论鳖臑]
鳖臑（四面体）的体积可以通过推理来阐明。按照阳马的体积，它是长方体的三分之一。而鳖臑又是阳马的一半。所以用长方体体积除以6就得到鳖臑体积。一般来说，一个长方体高1尺，宽和长各1尺，体积是1立方尺。阳马宽、长各1尺，高1尺，体积是1/3立方尺。鳖臑三条互相垂直的边各1尺，体积是1/6立方尺。
\end{commentarybox}

\end{paracol}

\vspace{0.5cm}
\longline
\vspace{0.3cm}

% ------------------------------------------------------------------------------
\bilingualsubsection{刍甍与刍童：楔形体}{The Chumeng and Chutong}
% ------------------------------------------------------------------------------

\begin{paracol}{2}

\begin{problembox}[\CJKfamily{rm}刍甍 · 原文]
\textbf{问：}今有刍甍，下广三丈，袤四丈，上袤二丈，无广，高一丈。问积几何？
\end{problembox}

\switchcolumn

\begin{problembox}[白话文：刍甍（屋脊楔形体）体积问题]
\textbf{题目：}现有一个刍甍，下底宽3丈，下底长4丈，上脊长2丈，上底无宽（为一条线），高1丈。问它的体积是多少？
\end{problembox}

\end{paracol}

\vspace{0.2cm}

\begin{paracol}{2}

\begin{answerbox}[\CJKfamily{rm}答 · 原文]
答曰：积五千尺。
\end{answerbox}

\switchcolumn

\begin{answerbox}[白话文：答案]
答案是：体积为5,000立方尺。
\end{answerbox}

\end{paracol}

\vspace{0.2cm}

\begin{paracol}{2}

\begin{methodbox}[\CJKfamily{rm}术 · 原文]
术曰：倍下袤，上袤从之，以广乘之，又以高乘之，六而一。
\end{methodbox}

\switchcolumn

\begin{methodbox}[白话文：刍甍体积算法]
算法说：将下底长度加倍，加上上脊长度，乘以底宽，再乘以高，除以6。
\end{methodbox}

\end{paracol}

\vspace{0.2cm}

\begin{paracol}{2}

\begin{commentarybox}[\CJKfamily{rm}刘徽注 · 原文]
推明义理者，旧说云：凡积刍甍有上下广曰童甍，谓其屋盖之苫也。是故甍之下广袤与童之上广袤等。正解方亭两边合之，即刍甍之形也。假令下广二尺，袤三尺，上袤一尺，无广，高一尺。其用棋也，中央堑堵二，两端阳马各二。倍下袤为六，上袤从之为七。以广二尺乘之，得一十四尺。以高乘之，得十四尺。六而一，即得积也。
\end{commentarybox}

\switchcolumn

\begin{commentarybox}[白话文：刘徽注释]
要阐明刍甍的体积原理，古人说过：凡计算刍甍，有上下底宽的就叫童甍，说的是屋顶的覆盖形状。所以甍的下底宽和长与童甍的上底宽和长相等。正好将方亭两边合在一起，就是刍甍的形状。假设下底宽2尺，长3尺，上脊长1尺，无上宽，高1尺。分解后，中央是两个堑堵，两端各两个阳马。倍下袤为6，加上上袤1得7。乘以底宽2尺，得14尺。再乘以高1尺，得14。除以6，就得到体积。
\end{commentarybox}

\end{paracol}

\vspace{0.3cm}

\begin{paracol}{2}

\begin{problembox}[\CJKfamily{rm}刍童 · 原文]
\textbf{问：}今有刍童，下广二丈，袤三丈，上广三丈，袤四丈，高三丈。问积几何？
\end{problembox}

\switchcolumn

\begin{problembox}[白话文：刍童（楔形截头体）体积问题]
\textbf{题目：}现有一个刍童，下底宽2丈，下底长3丈，上底宽3丈，上底长4丈，高3丈。问它的体积是多少？
\end{problembox}

\end{paracol}

\vspace{0.2cm}

\begin{paracol}{2}

\begin{answerbox}[\CJKfamily{rm}答 · 原文]
答曰：积一万六千五百尺。
\end{answerbox}

\switchcolumn

\begin{answerbox}[白话文：答案]
答案是：体积为16,500立方尺。
\end{answerbox}

\end{paracol}

\vspace{0.2cm}

\begin{paracol}{2}

\begin{methodbox}[\CJKfamily{rm}术 · 原文]
术曰：倍上袤，下袤从之；亦倍下袤，上袤从之。各以其广乘之，并，以高乘之，六而一。
\end{methodbox}

\switchcolumn

\begin{methodbox}[白话文：刍童体积算法]
算法说：将上底长度加倍，加上下底长度；再将下底长度加倍，加上上底长度。分别乘以各自的宽度，相加后乘以高度，除以6。
\end{methodbox}

\end{paracol}

\vspace{0.2cm}

\begin{paracol}{2}

\begin{commentarybox}[\CJKfamily{rm}刘徽注 · 原文]
按此术，假令刍童上广一尺，袤二尺；下广二尺，袤三尺；高一尺。其用棋也，中央立方一，两旁堑堵各一，四角阳马各一。倍上袤为四，下袤从之为七。以下广乘之，得一十四。倍下袤为六，上袤从之为八。以上广乘之，得八。并之，得二十二。以高乘之，得二十二。六而一，即积也。此术与刍甍曲池盘池冥谷皆同术。
\end{commentarybox}

\switchcolumn

\begin{commentarybox}[白话文：刘徽注释]
按照这个算法，假设刍童上底宽1尺，长2尺；下底宽2尺，长3尺；高1尺。分解后，中央是一个立方体，两旁各一个堑堵，四角各一个阳马。倍上袤为4，加上下袤3得7。以下底宽2乘之，得14。倍下袤为6，加上上袤2得8。以上底宽1乘之，得8。相加得22。乘以高1得22。除以6，就是体积。这个算法与刍甍、曲池、盘池、冥谷都是同一原理。
\end{commentarybox}

\end{paracol}

\vspace{0.5cm}
\longline
\vspace{0.3cm}

% ------------------------------------------------------------------------------
\bilingualsubsection{圆窖：圆形粮仓}{The Circular Granary}
% ------------------------------------------------------------------------------

\begin{paracol}{2}

\begin{problembox}[\CJKfamily{rm}圆窖 · 原文]
\textbf{问：}今有圆窖，下周五丈四尺，深一丈八尺。问受粟几何？
\end{problembox}

\switchcolumn

\begin{problembox}[白话文：圆窖（圆形粮仓）容量问题]
\textbf{题目：}现有一个圆形地窖，底面周长为5丈4尺，深1丈8尺。问它能装多少粟米？
\end{problembox}

\end{paracol}

\vspace{0.2cm}

\begin{paracol}{2}

\begin{answerbox}[\CJKfamily{rm}答 · 原文]
答曰：二千九百六十二斛二十七分斛之二十六。
\end{answerbox}

\switchcolumn

\begin{answerbox}[白话文：答案]
答案是：2,962又26/27斛。
\end{answerbox}

\end{paracol}

\vspace{0.2cm}

\begin{paracol}{2}

\begin{methodbox}[\CJKfamily{rm}术 · 原文]
术曰：置周自相乘，以高乘之，十二而一。以斛法一尺六寸二分除之，即斛数。
\end{methodbox}

\switchcolumn

\begin{methodbox}[白话文：圆窖容量算法]
算法说：将底面周长自乘，再乘以深度，除以12。再用一斛的容积（1立方尺6寸2分，即1.62立方尺）去除，就得到斛数。
\end{methodbox}

\end{paracol}

\vspace{0.2cm}

\begin{paracol}{2}

\begin{commentarybox}[\CJKfamily{rm}刘徽注 · 原文]
按此术，圆窖下周自乘，以高乘之，十二而一者，以圆周率三乘之，十二而一，即圆积也。此犹圆堡堭之积也。于徽术，当令下周自乘以高乘之，又以二十五乘之，三百一十四而一。以斛法除之，即斛数。
\end{commentarybox}

\switchcolumn

\begin{commentarybox}[白话文：刘徽注释]
按照这个算法，圆窖底面周长自乘，再乘以高度，除以12，这是用圆周率3来计算，十二分之一就是圆面积。这等同于圆堡堭的体积。按照刘徽的精确算法，应当将底面周长自乘，再乘以高度，再乘以25，除以314，然后用斛法去除，就得到斛数。
\end{commentarybox}

\end{paracol}

\newpage

% =============================================================================
\bilingualsection{卷六：均输}{Vol.\ 6: Junshu (Equitable Taxation)}
% =============================================================================

\begin{volumeheader}{\CJKfamily{rm}九章算術卷六 · 均输}
均输以御远近劳费\\[0.3em]
\small 白话文：均输章用于处理考虑远近、劳费等因素的公平分配
\end{volumeheader}

\vspace{0.5cm}

\begin{paracol}{2}
均输章研究加权分配问题，即根据不同地区的人口数、距目的地远近等因素公平分摊赋税和徭役。本章还包含著名的工作效率问题，如"五渠注池"。

\switchcolumn

均输章研究的是加权分配问题，即根据不同地区的人口数量、距离目的地的远近等因素，来公平合理地分摊赋税和劳役。本章还包含了许多著名的工作效率问题，其中最有名的是"五渠注池"（五条水渠同时注水）问题。

\end{paracol}

\vspace{0.5cm}

% ------------------------------------------------------------------------------
\bilingualsubsection{均输粟：四县分粮}{Opening Problem: Grain Distribution Among Four Counties}
% ------------------------------------------------------------------------------

\begin{paracol}{2}

\begin{problembox}[\CJKfamily{rm}均输粟 · 原文]
\textbf{问：}今有均输粟，甲县一万户，行道八日；乙县九千五百户，行道十日；丙县一万二千三百五十户，行道十三日；丁县一万二千二百户，行道二十日。各到输所。凡四县赋，当输二十五万斛，用车一万乘。欲以道里远近、户数多少，衰出之。问粟、车各几何？
\end{problembox}

\switchcolumn

\begin{problembox}[白话文：均输粮食问题]
\textbf{题目：}现在要公平地分配运粮任务：甲县有10,000户，到目的地的行程需要8天；乙县有9,500户，行程10天；丙县有12,350户，行程13天；丁县有12,200户，行程20天。各县都要把粮食运到指定地点。四个县共需缴纳25万斛粮食，动用1万辆车。要求按照路途远近和户数多少进行加权分摊。问各县应出多少粮食、出多少辆车？
\end{problembox}

\end{paracol}

\vspace{0.2cm}

\begin{paracol}{2}

\begin{answerbox}[\CJKfamily{rm}答 · 原文]
答曰：甲县粟八万三千一百斛，车三千三百二十四乘。乙县粟六万三千一百七十五斛，车二千五百二十七乘。丙县粟六万三千一百七十五斛，车二千五百二十七乘。丁县粟四万一千九百一十九斛，车一千六百二十二乘。
\end{answerbox}

\switchcolumn

\begin{answerbox}[白话文：答案]
答案是：甲县出粮83,100斛，出车3,324辆。乙县出粮63,175斛，出车2,527辆。丙县出粮63,175斛，出车2,527辆。丁县出粮41,919斛，出车1,622辆。
\end{answerbox}

\end{paracol}

\vspace{0.2cm}

\begin{paracol}{2}

\begin{methodbox}[\CJKfamily{rm}术 · 原文]
术曰：令县户数，各如其本行道日数而一，以为衰。甲衰一百二十五，乙、丙衰各九十五，丁衰六十一。副并为法。以赋粟、车数，未并者各自为实。实如法得一。按此均输，犹均运也。令户率出车，以行道日数为均，发粟为输。
\end{methodbox}

\switchcolumn

\begin{methodbox}[白话文：均输加权分配算法]
算法说：令各县的户数分别除以各自的天数，作为衰减率（衰）。甲县的衰为125，乙县和丙县的衰各为95，丁县的衰为61。将所有衰相加作为除数（法）。用赋粟总数和车辆总数分别乘以各县未相加前的衰作为被除数（实）。被除数除以除数就得到各县的份额。这种均输法，就是均匀运输的意思。让每户按比例出车，以行程天数作为均等的依据，发出粮食就是运输。
\end{methodbox}

\end{paracol}

\vspace{0.2cm}

\begin{paracol}{2}

\begin{commentarybox}[\CJKfamily{rm}刘徽注 · 原文]
此户以率当各计车之钱分也。案此二句舛误，当云：求此率以户当各计车之钱分也。淳风等按：县户有多少之差，行道有远近之异。欲其均等，故令户数各以行道日数约之，为衰。行道多者，其户少；行道少者，其户多。故各令约户为衰。并户为法者，以户数多寡为差，行道远近为衰，以此均之，则劳逸齐等。
\end{commentarybox}

\switchcolumn

\begin{commentarybox}[白话文：刘徽与李淳风注释]
这里是要按比率计算每户应分摊的车辆和费用。这两句文字有误，应当是：求此率，以户当各计车之钱分也。李淳风等按：各县户数有多有少，行程有远有近。要使劳逸均等，所以让各县的户数分别除以行程天数，作为衰减率。行程天数多的县，分摊到的户数就少；行程天数少的县，分摊到的户数就多。所以让各县以户数除以天数作为衰。将各县衰相加作为法，以户数多少为差、行程远近为衰，以此来均分，就能使劳逸齐等。
\end{commentarybox}

\end{paracol}

\vspace{0.5cm}
\longline
\vspace{0.3cm}

% ------------------------------------------------------------------------------
\bilingualsubsection{程人功：工作效率问题}{Workers with Different Rates}
% ------------------------------------------------------------------------------

\begin{paracol}{2}

\begin{problembox}[\CJKfamily{rm}程人功 · 原文]
\textbf{问：}今有程人功，甲七日成功了五分之一；乙五日成功了三分之一；丙四日成功了二分之一。问三人合作，几何日成功？
\end{problembox}

\switchcolumn

\begin{problembox}[白话文：工作效率问题]
\textbf{题目：}现有三名工人：甲7天完成了全部工作的五分之一；乙5天完成了全部工作的三分之一；丙4天完成了全部工作的二分之一。问三人一起合作，需要多少天才能完成全部工作？
\end{problembox}

\end{paracol}

\vspace{0.2cm}

\begin{paracol}{2}

\begin{answerbox}[\CJKfamily{rm}答 · 原文]
答曰：二十三日六十三分日之二十六。
\end{answerbox}

\switchcolumn

\begin{answerbox}[白话文：答案]
答案是：23又26/63天。（约23.41天）
\end{answerbox}

\end{paracol}

\vspace{0.2cm}

\begin{paracol}{2}

\begin{methodbox}[\CJKfamily{rm}术 · 原文]
术曰：置甲七日的五分之一，乙五日的三分之一，丙四日的二分之一，各为率。以一日为法，以所为实。实如法而一，各得一日之功。并一日之功，为法。以一日为实。实如法而一，即得日数。
\end{methodbox}

\switchcolumn

\begin{methodbox}[白话文：工作效率算法]
算法说：列出甲7天完成五分之一、乙5天完成三分之一、丙4天完成二分之一，各自作为比率。以一天为标准（法），各自完成的量为被除数（实）。相除得到每人一天的工作量。将三人一天的工作量相加作为除数（法），以一天为单位作为被除数（实）。被除数除以除数，就得到完成全部工作需要的天数。
\end{methodbox}

\end{paracol}

\vspace{0.5cm}
\longline
\vspace{0.3cm}

% ------------------------------------------------------------------------------
\bilingualsubsection{五渠注池：五渠灌水问题}{Five Channels Filling a Pool}
% ------------------------------------------------------------------------------

这是《九章算术》中最著名的问题之一，作为卷六的最后一题出现。

\begin{paracol}{2}

\begin{problembox}[\CJKfamily{rm}五渠注池 · 原文]
\textbf{问：}今有池，五渠注之。其一渠开之，少半日一满；次，一日一满；次，二日半一满；次，三日一满；次，五日一满。今并决之，问几何日满池？
\end{problembox}

\switchcolumn

\begin{problembox}[白话文：五渠注池问题]
\textbf{题目：}现有一个水池，有五条水渠向它注水。第一条渠单独开，1/3天就注满；第二条渠单独开，1天注满；第三条渠单独开，2.5天注满；第四条渠单独开，3天注满；第五条渠单独开，5天注满。现在五条渠同时打开，问多少天可以注满水池？
\end{problembox}

\end{paracol}

\vspace{0.2cm}

\begin{paracol}{2}

\begin{answerbox}[\CJKfamily{rm}答 · 原文]
答曰：七十四分日之十五。
\end{answerbox}

\switchcolumn

\begin{answerbox}[白话文：答案]
答案是：15/74天。（约0.203天，约4.86小时）
\end{answerbox}

\end{paracol}

\vspace{0.2cm}

\begin{paracol}{2}

\begin{methodbox}[\CJKfamily{rm}术 · 原文]
术曰：各置渠一日满池之数，并，以为法。以一为实。实如法而一，即得日数。
\end{methodbox}

\switchcolumn

\begin{methodbox}[白话文：五渠注水算法]
算法说：分别列出每条水渠一天能注满水池的份数（即各渠的日注水效率），相加作为除数（法）。以注满一池作为被除数（实）。被除数除以除数，就得到注满水池需要的天数。
\end{methodbox}

\end{paracol}

\vspace{0.2cm}

\begin{paracol}{2}

\begin{commentarybox}[\CJKfamily{rm}刘徽注 · 原文]
按此术，一渠少半日满者，是一日三满也。次一日一满者，是一日一满也。次二日半满者，是一日五分满之二也。次三日一满者，是一日三分满之一也。次五日一满者，是一日五分满之一也。并之，得四满十五分满之十四也。此为法。以一满为实。实如法而一，即得满池之日数也。
\end{commentarybox}

\switchcolumn

\begin{commentarybox}[白话文：刘徽注释]
按照这个算法，第一条渠1/3天就注满，意味着一天能注满3池。第二条渠一天注满1池。第三条渠2.5天注满，意味着一天能注满2/5池。第四条渠三天注满，意味着一天能注满1/3池。第五条渠五天注满，意味着一天能注满1/5池。将这些相加：$3 + 1 + \frac{2}{5} + \frac{1}{3} + \frac{1}{5} = \frac{74}{15}$池/天。这就是除数（法）。以注满一池作为被除数（实）。被除数除以除数，就得到注满水池需要的天数。
\end{commentarybox}

\end{paracol}

\vspace{0.5cm}
\longline
\vspace{0.3cm}

% ------------------------------------------------------------------------------
\bilingualsubsection{兔犬问题：追及问题}{Hare and Dog Pursuit}
% ------------------------------------------------------------------------------

\begin{paracol}{2}

\begin{problembox}[\CJKfamily{rm}兔犬 · 原文]
\textbf{问：}今有兔先走一百步，犬追之二百五十步，不及三十步而止。问犬不止，复行几何步及之？
\end{problembox}

\switchcolumn

\begin{problembox}[白话文：兔犬追及问题]
\textbf{题目：}一只兔子先跑了100步，一条狗追了250步后，还差30步没追上就停下来了。问如果狗不停止，继续追，还需要跑多少步才能追上兔子？
\end{problembox}

\end{paracol}

\vspace{0.2cm}

\begin{paracol}{2}

\begin{answerbox}[\CJKfamily{rm}答 · 原文]
答曰：一百七步七分步之一。
\end{answerbox}

\switchcolumn

\begin{answerbox}[白话文：答案]
答案是：107又1/7步。
\end{answerbox}

\end{paracol}

\vspace{0.2cm}

\begin{paracol}{2}

\begin{methodbox}[\CJKfamily{rm}术 · 原文]
术曰：置犬先行步数，以不及步数减之，余为法。以不及步数乘兔先走步数，为实。实如法得一步。
\end{methodbox}

\switchcolumn

\begin{methodbox}[白话文：追及问题算法]
算法说：列出狗已经追的步数，用尚未追上的步数去减，余数作为除数（法）。用尚未追上的步数乘以兔子先跑的步数，作为被除数（实）。被除数除以除数，就得到还需要追的步数。
\end{methodbox}

\end{paracol}

\vspace{0.5cm}
\longline
\vspace{0.3cm}

% ------------------------------------------------------------------------------
\bilingualsubsection{金银问题：方程术}{Gold and Silver Weights}
% ------------------------------------------------------------------------------

\begin{paracol}{2}

\begin{problembox}[\CJKfamily{rm}金银 · 原文]
\textbf{问：}今有金九枚，银一十一枚，称之适等。交易其一，金轻十三两。问金、银各一枚重几何？
\end{problembox}

\switchcolumn

\begin{problembox}[白话文：金银重量问题]
\textbf{题目：}现有9枚金币和11枚银币，称量后重量恰好相等。交换其中各一枚后，金币这边轻了13两。问每枚金币和每枚银币各重多少？
\end{problembox}

\end{paracol}

\vspace{0.2cm}

\begin{paracol}{2}

\begin{answerbox}[\CJKfamily{rm}答 · 原文]
答曰：金重二斤三两一十八铢，银重一斤十三两六铢。
\end{answerbox}

\switchcolumn

\begin{answerbox}[白话文：答案]
答案是：每枚金币重2斤3两18铢，每枚银币重1斤13两6铢。
\end{answerbox}

\end{paracol}

\vspace{0.2cm}

\begin{paracol}{2}

\begin{methodbox}[\CJKfamily{rm}术 · 原文]
术曰：假令金重三斤，银重二斤三分斤之二。不足四两，令之如此。盈不足术为之。
\end{methodbox}

\switchcolumn

\begin{methodbox}[白话文：金银重量算法]
算法说：假设每枚金币重3斤，每枚银币重2又2/3斤。发现这样算会差4两，就按这个假设来做。用盈不足术来求解。
\end{methodbox}

\end{paracol}

\newpage

% =============================================================================
\section{数学内容总结 · Summary of Mathematical Content}
% =============================================================================

\begin{paracol}{2}

\textbf{卷四：少广 (Shaoguang)}

\begin{itemize}[leftmargin=1.5em]
  \item \textbf{开方术}：求 $\sqrt{N}$ 的迭代算法，本质上等同于现代开平方的长除法
  \item \textbf{开立方术}：扩展到 $\sqrt[3]{N}$
  \item \textbf{球体积}：古法 $V = \frac{9}{16}d^3$；刘徽的批评与牟合方盖构造
  \item 带分数运算及其几何解释
\end{itemize}

\switchcolumn

\textbf{白话文总结：}

\begin{itemize}[leftmargin=1.5em]
  \item \textbf{开平方}：通过迭代算法求一个数的平方根，原理与现代的长除法开平方相同
  \item \textbf{开立方}：将开平方的方法推广到求立方根
  \item \textbf{球体积}：古法采用 $V = \frac{9}{16}d^3$ 的近似公式；刘徽发现其错误，提出牟合方盖方法，但未完成最终求解
  \item 分数与几何图形相结合的计算方法
\end{itemize}

\end{paracol}

\vspace{0.5cm}

\begin{paracol}{2}

\textbf{卷五：商功 (Shanggong)}

\begin{itemize}[leftmargin=1.5em]
  \item \textbf{土方体积}：城、垣、堤、沟、渠、堑
  \item \textbf{仓窖体积}：仓、圆窖、圆囤
  \item \textbf{几何体}：立方、堑堵、阳马、鳖臑、刍甍、刍童、曲池、盘池、冥谷
  \item 刘徽的\textbf{几何分割证明}
\end{itemize}

\switchcolumn

\textbf{白话文总结：}

\begin{itemize}[leftmargin=1.5em]
  \item \textbf{工程土方量}：城墙、普通墙、堤坝、水沟、水渠、堑壕的体积计算
  \item \textbf{粮仓容量}：方形粮仓、圆形地窖、圆形粮囤的容量计算
  \item \textbf{几何体体积}：立方体、楔形、阳马（直角方锥）、鳖臑（四面体）、刍甍（屋脊形）、刍童（截头楔形体）、曲池、盘池、冥谷等
  \item 刘徽运用几何分割和拼补法给出的严格体积证明
\end{itemize}

\end{paracol}

\vspace{0.5cm}

\begin{paracol}{2}

\textbf{卷六：均输 (Junshu)}

\begin{itemize}[leftmargin=1.5em]
  \item \textbf{加权分配}：按人口$\times$距离分摊赋税
  \item \textbf{工作效率}：多人合作问题
  \item \textbf{追及问题}：相对速度（兔犬问题）
  \item \textbf{方程术}：金银重量、粮食交换
  \item 著名的\textbf{五渠注池}问题
\end{itemize}

\switchcolumn

\textbf{白话文总结：}

\begin{itemize}[leftmargin=1.5em]
  \item \textbf{加权分配}：根据各县人口数量和距离远近的乘积来公平分摊赋税和运输任务
  \item \textbf{工作效率}：多名工人以不同效率合作完成工作的问题
  \item \textbf{追及问题}：利用相对速度求解追赶问题（如兔犬问题）
  \item \textbf{方程术}：通过设立方程求解金银重量、粮食交换等问题
  \item 著名的\textbf{五条水渠同时注水}的效率问题
\end{itemize}

\end{paracol}

\vspace{1cm}

\begin{center}
\shortline
\vspace{0.5cm}

\textit{「九章之术，凡数之本。」}
\vspace{0.3em}

--- 刘徽《九章算术注序》

\vspace{0.3cm}

\textit{白话文："《九章算术》的方法，是一切数学计算的根源。"}

\vspace{0.5cm}
\shortline
\end{center}

\vspace{1cm}

\begin{center}
\textbf{\CJKfamily{rm}九章算術卷四至卷六 完}\\[0.5em]
\textit{End of Jiuzhang Suanshu, Volumes 4--6}
\end{center}

\end{document}
